\documentclass[12pt, a4paper]{article}

% -- Packages -----------------------------------------------------------------
\usepackage[margin=1in]{geometry}
\usepackage[T1]{fontenc}
\usepackage[utf8]{inputenc}
\usepackage{lmodern}
\usepackage{microtype}
\usepackage{xcolor}
\usepackage{titlesec}
\usepackage{booktabs}
\usepackage{tabularx}
\usepackage{longtable}
\usepackage{array}
\usepackage{colortbl}
\usepackage{enumitem}
\usepackage{listings}
\usepackage{hyperref}
\usepackage{fancyhdr}
\usepackage{graphicx}
\usepackage{tcolorbox}
\usepackage{setspace}
\usepackage{mdframed}
\usepackage{amsmath}
\tcbuselibrary{skins, breakable}

% -- Colors -------------------------------------------------------------------
\definecolor{primaryblue}{HTML}{1A73E8}
\definecolor{darkblue}{HTML}{0D47A1}
\definecolor{lightgray}{HTML}{F5F5F5}
\definecolor{codebg}{HTML}{1E1E2E}
\definecolor{codetext}{HTML}{CDD6F4}
\definecolor{accent}{HTML}{E53935}
\definecolor{mysqlorange}{HTML}{E8710A}
\definecolor{tableheader}{HTML}{1A73E8}
\definecolor{successgreen}{HTML}{1B8A3E}
\definecolor{lightgreen}{HTML}{F0FFF4}

% -- Hyperref Setup -----------------------------------------------------------
\hypersetup{
    colorlinks=true,
    linkcolor=primaryblue,
    urlcolor=primaryblue,
    pdftitle={GitDB: Stack and Feature Implementation Guide},
    pdfauthor={Group Project - DBMS}
}

% -- Section Formatting --------------------------------------------------------
\titleformat{\section}
  {\large\bfseries\color{darkblue}}
  {\thesection.}{0.5em}{}[\color{primaryblue}\titlerule]

\titleformat{\subsection}
  {\normalsize\bfseries\color{primaryblue}}
  {\thesubsection.}{0.5em}{}

\titleformat{\subsubsection}
  {\normalsize\bfseries\color{primaryblue!80}}
  {\thesubsubsection.}{0.5em}{}

\titlespacing*{\section}{0pt}{1.2em}{0.6em}
\titlespacing*{\subsection}{0pt}{0.8em}{0.4em}
\titlespacing*{\subsubsection}{0pt}{0.6em}{0.3em}
\setlength{\LTpre}{0pt}
\setlength{\LTpost}{0.5em}

% -- Code Listing Style -------------------------------------------------------
\lstset{
    backgroundcolor=\color{codebg},
    basicstyle=\ttfamily\small\color{codetext},
    keywordstyle=\color{cyan}\bfseries,
    commentstyle=\color{gray}\itshape,
    stringstyle=\color{yellow},
    breaklines=true,
    frame=none,
    xleftmargin=1em,
    xrightmargin=1em,
    aboveskip=0.5em,
    belowskip=0.5em,
    showstringspaces=false,
    numbers=none,
}

% -- Custom Boxes -------------------------------------------------------------
\tcbset{
    commandbox/.style={
        colback=codebg,
        colframe=primaryblue,
        coltext=codetext,
        fonttitle=\bfseries\color{white},
        attach boxed title to top left={yshift=-2mm, xshift=4mm},
        boxed title style={colback=primaryblue, colframe=primaryblue},
        arc=4pt,
        boxrule=0.8pt,
        left=6pt, right=6pt, top=8pt, bottom=6pt,
        breakable
    },
    infobox/.style={
        colback=lightgray,
        colframe=primaryblue,
        arc=4pt,
        boxrule=0.8pt,
        left=8pt, right=8pt, top=6pt, bottom=6pt
    },
    warningbox/.style={
        colback=mysqlorange!8,
        colframe=mysqlorange,
        arc=4pt,
        boxrule=0.8pt,
        left=8pt, right=8pt, top=6pt, bottom=6pt
    },
    fixbox/.style={
        colback=lightgreen,
        colframe=successgreen,
        arc=4pt,
        boxrule=0.8pt,
        left=8pt, right=8pt, top=6pt, bottom=6pt
    }
}

% -- Header / Footer ----------------------------------------------------------
\pagestyle{fancy}
\fancyhf{}
\fancyhead[L]{\small\textcolor{primaryblue}{\textbf{GitDB}} --- Stack and Features}
\fancyhead[R]{\small\textcolor{gray}{DBMS Mini Project}}
\fancyfoot[C]{\small\textcolor{gray}{\thepage}}
\renewcommand{\headrulewidth}{0.4pt}
\renewcommand{\headrule}{\hbox to\headwidth{\color{primaryblue}\leaders\hrule height \headrulewidth\hfill}}

% -- Document -----------------------------------------------------------------
\begin{document}

% -- Title Page ---------------------------------------------------------------
\begin{titlepage}
    \centering
    \vspace*{1.5cm}

    {\Huge\bfseries\color{darkblue} GitDB}\\[0.4em]
    {\Large\color{primaryblue} Technology Stack and Feature Implementation Guide}\\[0.2em]
    {\small\color{mysqlorange}\textit{What is used, what is custom-built, and how it works}}\\[0.2em]
    {\small\color{gray}\rule{8cm}{0.5pt}}

    \vspace{1.5cm}

    \begin{tcolorbox}[infobox, width=0.88\textwidth]
        \centering
        {\normalsize\textit{This document explains the full GitDB project implementation from source code:\\
        backend stack, frontend stack, database design, all major features, and\\
        whether each feature is implemented through external packages or custom Python/JavaScript code.}}
    \end{tcolorbox}

    \vspace{1.5cm}

    {\large\textbf{Institution:} PICT, Pune}\\[0.5em]
    {\large\textbf{Subject:} DISL}\\[0.5em]
    {\large\textbf{Project Type:} Mini Project}\\[0.5em]
    {\large\textbf{Class:} SY-11 \quad|\quad \textbf{Lab Batch:} F-11}\\[0.5em]
    {\large\textbf{Group Size:} 4 Members}\\[1em]

    \vspace{0.5cm}

    \textbf{Team Members:}\\
    \vspace{0.5em}
    \renewcommand{\arraystretch}{1.3}
    \begin{tabular}{>{\centering\arraybackslash}p{1cm} >{\arraybackslash}p{5cm} >{\centering\arraybackslash}p{2.5cm}}
    \toprule
    \textbf{Sr.} & \textbf{Name} & \textbf{Roll No.} \\
    \midrule
    1 & Aditya Malode     & 23325 \\
    2 & Rohan Natekar     & 23330 \\
    3 & Aarush Oak        & 23332 \\
    4 & Prathamesh Kinagi & 23340 \\
    \bottomrule
    \end{tabular}

    \vspace{2.5cm}
    {\small\color{gray} Academic Year 2025--26}
    \vfill
    {\small\color{gray} Compiled on \today}
\end{titlepage}

\tableofcontents
\newpage

\section{Project Summary}

GitDB is a Git-like version control system for MySQL databases. It tracks full
\textbf{state snapshots} (schema plus data) instead of migration instructions.
Each commit stores a full snapshot inside \texttt{gitdb\_meta}, computes a
SHA-256 commit hash, and links to a parent hash to form a linear Merkle chain.

The project is implemented as a full-stack system:
\begin{itemize}[leftmargin=1.5em]
    \item \textbf{Core engine (Python):} snapshotting, hashing, diffing, checkout
    \item \textbf{CLI (Python):} user-facing commands similar to Git
    \item \textbf{REST API (Flask):} UI-friendly endpoints
    \item \textbf{Frontend (React + Vite):} commit graph, diff viewer, schema browser
    \item \textbf{Metadata DB (MySQL \texttt{gitdb\_meta}):} persistent versioning objects
    \item \textbf{Automated tests (pytest):} unit plus integration validation
\end{itemize}

\begin{tcolorbox}[infobox]
\textbf{Design principle:} keep the version-control layer separate from the target
application database. GitDB metadata is stored in a dedicated \texttt{gitdb\_meta}
database, similar to how Git stores metadata in a dedicated \texttt{.git/} directory.
\end{tcolorbox}

\section{Technology Stack}

\subsection{Backend and Database Stack}

\renewcommand{\arraystretch}{1.3}
\begin{longtable}{@{}>{\bfseries}p{2.7cm} >{\ttfamily}p{3.3cm} p{9.3cm}@{}}
\toprule
Layer & Technology & Purpose \\
\midrule
Language & Python 3.10+ & Primary implementation language for engine, CLI, and API \\
Database driver & mysql-connector-python & Connect to MySQL and execute SQL \\
SQL parser & sqlglot & Parse \texttt{CREATE TABLE} DDL into AST for structural schema diffing \\
CLI framework & click & Command definitions, options, arguments, command groups \\
CLI rendering & rich & Colored tables, panels, and readable terminal output \\
Web API & flask + flask-cors & REST service and browser-safe cross-origin access \\
Password hashing & argon2-cffi & Secure hash for GitDB user passwords \\
Secret storage & keyring & Store target DB passwords in OS keychain instead of plaintext config \\
Cryptography & hashlib (standard library) & SHA-256 commit hash generation \\
Testing & pytest & Unit and integration test framework \\
\bottomrule
\end{longtable}

\subsection{Frontend Stack}

\renewcommand{\arraystretch}{1.3}
\begin{longtable}{@{}>{\bfseries}p{2.7cm} >{\ttfamily}p{3.3cm} p{9.3cm}@{}}
\toprule
Layer & Technology & Purpose \\
\midrule
UI runtime & React 18 & Component-based frontend \\
Build tool & Vite 5 & Fast development server and production build pipeline \\
Styling & Tailwind CSS 4 + custom CSS variables & Utility-first styling plus custom themed design system \\
Commit graph & @xyflow/react + dagre & Interactive graph nodes/edges plus automatic layout \\
Visual diff & @git-diff-view/react + @git-diff-view/file & Side-by-side schema diff rendering \\
Browser API client & fetch (custom wrapper) & API requests to Flask endpoints \\
\bottomrule
\end{longtable}

\subsection{Persistence Model}

GitDB metadata database schema is implemented in \texttt{db/schema.sql} using
four core entities:
\begin{itemize}[leftmargin=1.5em]
    \item \texttt{user}: authentication and commit authorship
    \item \texttt{repository}: target DB config and current HEAD hash
    \item \texttt{commits}: linear commit chain with parent references
    \item \texttt{snapshots}: per-table serialized schema and row data
\end{itemize}

\section{Codebase Structure and Responsibilities}

\begin{center}
\renewcommand{\arraystretch}{1.25}
\begin{tabularx}{\textwidth}{>{\ttfamily}p{3.3cm} X}
\toprule
Path & Responsibility \\
\midrule
engine/snapshot.py & Capture schema and row snapshots from target MySQL DB \\
engine/hash.py & Deterministic commit hash computation \\
engine/diff.py & Schema and data diff generation as SQL patch statements \\
engine/checkout.py & Two-phase checkout executor with recovery logic \\
engine/db.py & Config loading, keyring access, target/meta DB connection helpers \\
engine/errors.py & Shared custom exception types \\
cli/main.py & All CLI commands: register/init/commit/log/diff/checkout/status \\
api/app.py & REST endpoints for commits/diff/snapshot/checkout/status \\
ui/src/App.jsx & Main UI shell and page orchestration \\
ui/src/components/*.jsx,*.js & Commit node rendering and graph layout logic \\
tests/unit & Fast isolated unit tests with mocks \\
tests/integration & Real MySQL integration scenarios \\
\bottomrule
\end{tabularx}
\end{center}

\section{GitDB Data Structures}

This section describes the core data structures used internally by GitDB.
These structures are the implementation backbone of commit creation, diff
generation, and checkout execution.

\subsection{SQL Metadata Database Structure (\texttt{gitdb\_meta})}

GitDB stores all version-control metadata in a dedicated MySQL database named
\texttt{gitdb\_meta}. This is the primary database-side data structure of the
system, and it is defined in \texttt{db/schema.sql}.

\begin{center}
\renewcommand{\arraystretch}{1.3}
\begin{tabularx}{\textwidth}{>{\bfseries}p{2.7cm} >{\ttfamily}p{4.0cm} X}
\toprule
Table & Key fields & Role in GitDB \\
\midrule
user & \texttt{user\_id} (PK), \texttt{username} (UNIQUE), \texttt{email} (UNIQUE) & Stores GitDB login identity and commit authorship metadata \\
repository & \texttt{repo\_id} (PK), \texttt{user\_id} (FK), \texttt{current\_hash} (FK) & Stores target DB connection metadata and HEAD pointer \\
commits & \texttt{hash} (PK), \texttt{repo\_id} (FK), \texttt{parent\_hash} (self-FK, UNIQUE) & Stores linear commit chain and commit message/timestamp/author \\
snapshots & \texttt{id} (PK), \texttt{commit\_hash} (FK), \texttt{(commit\_hash, table\_name)} (UNIQUE) & Stores per-table schema/data snapshot payload for each commit \\
\bottomrule
\end{tabularx}
\end{center}

\begin{tcolorbox}[infobox]
\textbf{Key design point:}
\texttt{repository.current\_hash} acts as the HEAD reference, and
\texttt{commits.parent\_hash} creates a strict linear history structure.
\end{tcolorbox}

\begin{tcolorbox}[commandbox, title=Relationship Structure (Conceptual)]
\begin{lstlisting}
user (1) -------- (N) repository (1) -------- (N) commits (1) -------- (N) snapshots
                           |                        |
                           |                        +-- parent_hash -> commits.hash (self link)
                           +-- current_hash -> commits.hash (HEAD pointer)
\end{lstlisting}
\end{tcolorbox}

\subsection{In-Memory Core Structures (Python)}

\begin{center}
\renewcommand{\arraystretch}{1.3}
\begin{tabularx}{\textwidth}{>{\bfseries}p{3.0cm} >{\ttfamily}p{3.1cm} X}
\toprule
Structure & Declared in & Purpose \\
\midrule
TableSnapshot & engine/snapshot.py & Represents one table state: DDL, rows, row count, and PK metadata \\
RepoConfig & engine/db.py & Runtime repository connection model loaded from \texttt{.gitdb/config.json} \\
DiffResult & engine/diff.py & Holds generated schema SQL, data SQL, and warning messages \\
Snapshot map & engine/snapshot.py & Dictionary \texttt{dict[str, TableSnapshot]} for full database state at one commit \\
PK index maps & engine/diff.py & Row lookup maps keyed by PK tuple for efficient row-level diff \\
\bottomrule
\end{tabularx}
\end{center}

\begin{tcolorbox}[commandbox, title=Key Python Data Structures (Conceptual Form)]
\begin{lstlisting}
TableSnapshot:
    table_name: str
    ddl_json: dict
    rows_json: list[dict]
    row_count: int
    pk_columns: list[str]

RepoConfig:
    repo_id: int
    host: str
    port: int
    db_user: str
    db_name: str

DiffResult:
    schema_sql: list[str]
    data_sql: list[str]
    warnings: list[str]
\end{lstlisting}
\end{tcolorbox}

\subsection{Commit History Data Structure}

GitDB commit history is a \textbf{linear linked structure} (no branching).
At persistence level, this is represented by:
\begin{itemize}[leftmargin=1.5em]
        \item \texttt{commits.hash}: node identity
        \item \texttt{commits.parent\_hash}: pointer to previous node
        \item \texttt{repository.current\_hash}: HEAD pointer to latest checked-out node
        \item unique \texttt{parent\_hash}: ensures at most one child per commit in this design
\end{itemize}

\begin{tcolorbox}[infobox]
Conceptually, this behaves like a singly linked list with cryptographic node IDs:
HEAD $\rightarrow C_n \rightarrow C_{n-1} \rightarrow \dots \rightarrow C_0$.
\end{tcolorbox}

\subsection{Snapshot Storage Structure}

Each commit stores a full database snapshot as a set of per-table records.
The persistence pattern is:
\begin{itemize}[leftmargin=1.5em]
        \item one row in \texttt{commits}
        \item many rows in \texttt{snapshots} (exactly one per table name)
        \item \texttt{ddl\_json} stores schema metadata and raw DDL
        \item \texttt{rows\_json} stores deterministic row arrays ordered by PK columns
\end{itemize}

\begin{tcolorbox}[commandbox, title=Snapshot Container Structure]
\begin{lstlisting}
snapshot: dict[str, TableSnapshot] = {
    "users": TableSnapshot(...),
    "orders": TableSnapshot(...),
    "payments": TableSnapshot(...),
}
\end{lstlisting}
\end{tcolorbox}

\subsection{Diff Engine Working Structures}

During diff generation, GitDB transforms row arrays into hash-map style indexes
to support efficient set operations:
\begin{itemize}[leftmargin=1.5em]
        \item \texttt{old\_idx = \{pk\_tuple(row): row\}} and \texttt{new\_idx = \{pk\_tuple(row): row\}}
        \item \texttt{added = new\_keys - old\_keys}
        \item \texttt{removed = old\_keys - new\_keys}
        \item \texttt{common = old\_keys \& new\_keys}
\end{itemize}

These structures drive deterministic emission of \texttt{INSERT},
\texttt{DELETE}, and \texttt{UPDATE} SQL statements.

\subsection{API Payload Structures}

GitDB API responses follow stable JSON object/array structures, for example:
\begin{itemize}[leftmargin=1.5em]
        \item \texttt{/commits}: list of commit objects with hash, parent, message, timestamp, author
        \item \texttt{/diff}: object with \texttt{schema\_sql}, \texttt{data\_sql}, \texttt{warnings}
        \item \texttt{/snapshot}: list of table snapshot summaries (table name, row count, ddl)
        \item \texttt{/status}: object with table additions/deletions/modifications and row deltas
\end{itemize}

\begin{tcolorbox}[fixbox]
\textbf{Why this matters:}
Clearly defined data structures make GitDB testable and maintainable. Most unit
tests directly validate behavior of these structures (hash determinism, snapshot
shape, diff SQL generation, checkout state transitions, and API contract shape).
\end{tcolorbox}

\section{Feature-Wise Implementation Mapping}

\subsection{Quick Mapping Table}

\begin{center}
\renewcommand{\arraystretch}{1.35}
\begin{tabularx}{\textwidth}{>{\bfseries}p{3.0cm} >{\bfseries}p{3.2cm} >{\bfseries}p{3.2cm} X}
\toprule
Feature & Package-assisted part & Custom code part & Main implementation \\
\midrule
User registration & \texttt{argon2-cffi} & Validation and DB insert flow & \texttt{cli/main.py} \\
Repo init and config & \texttt{click}, \texttt{keyring} & Metadata setup and config file writing & \texttt{cli/main.py}, \texttt{engine/db.py} \\
Snapshot capture & \texttt{mysql-connector-python} & Introspection, PK handling, row-limit guard & \texttt{engine/snapshot.py} \\
Commit identity & \texttt{hashlib} & Payload construction and deterministic serialization path & \texttt{engine/hash.py} \\
Schema diff & \texttt{sqlglot} AST parser & Added/removed/changed column SQL generation & \texttt{engine/diff.py} \\
Data diff & None (algorithm is custom) & PK-indexed \texttt{INSERT}/\texttt{UPDATE}/\texttt{DELETE} generation & \texttt{engine/diff.py} \\
Checkout safety & MySQL transaction primitives & Two-phase execution, rollback, recovery file path & \texttt{engine/checkout.py} \\
CLI UX & \texttt{click}, \texttt{rich} & Command semantics and orchestration & \texttt{cli/main.py} \\
REST API & \texttt{flask}, \texttt{flask-cors} & Endpoint logic and JSON formatting & \texttt{api/app.py} \\
UI graph and diff & \texttt{@xyflow/react}, \texttt{dagre}, \texttt{@git-diff-view/react} & App flow, API integration, page state & \texttt{ui/src/App.jsx} \\
\bottomrule
\end{tabularx}
\end{center}

\subsection{Feature 1: User Registration and Authentication Metadata}

\textbf{What it does:}
Creates GitDB users in \texttt{gitdb\_meta.user} and stores secure password hashes.

\textbf{How it is implemented:}
\begin{itemize}[leftmargin=1.5em]
    \item Custom CLI flow in \texttt{register} command (inputs, DB setup, insert)
    \item Password hashing delegated to \texttt{argon2-cffi}
    \item User lookup by username used later for commit authorship
\end{itemize}

\textbf{Why package is used:}
Argon2 is a secure password hashing standard. Writing this manually is unsafe,
so a proven cryptographic package is used.

\subsection{Feature 2: Repository Initialization and Secure Credential Handling}

\textbf{What it does:}
Initializes metadata for one target DB, stores non-secret config in
\texttt{.gitdb/config.json}, and stores the real DB password in OS keyring.

\textbf{How it is implemented:}
\begin{itemize}[leftmargin=1.5em]
    \item Custom init command creates \texttt{repository} row and writes config file
    \item \texttt{engine/db.py} provides \texttt{load\_repo\_config()} and connection helpers
    \item \texttt{keyring.set\_password()} and \texttt{get\_password()} provide secret storage
\end{itemize}

\textbf{Implementation ownership:}
Control flow and repository metadata management are custom Python code. Secret vault
behavior is delegated to the OS keyring package.

\subsection{Feature 3: Snapshot Capture (Schema + Data)}

\textbf{What it does:}
Captures the full table-level state of the target DB for each commit.

\textbf{How it is implemented:}
\begin{itemize}[leftmargin=1.5em]
    \item Queries \texttt{information\_schema.TABLES} for base table list
    \item Queries \texttt{information\_schema.COLUMNS} and \texttt{SHOW CREATE TABLE}
    \item Detects primary-key columns via \texttt{information\_schema.KEY\_COLUMN\_USAGE}
    \item Reads rows using \texttt{SELECT * ORDER BY <pk columns>} for deterministic ordering
    \item Enforces hard guard \texttt{row\_limit = 10000} per table
    \item For no-PK tables: captures DDL but skips row data diff payload
\end{itemize}

\textbf{Package vs custom:}
Database access is package-driven (MySQL connector), but snapshot data model,
introspection sequence, row-limit policy, and no-PK behavior are custom logic.

\subsection{Feature 4: Commit Hashing and Linear Merkle Chain}

\textbf{What it does:}
Computes immutable commit identifiers and links each commit to a parent.

\textbf{How it is implemented:}
\begin{itemize}[leftmargin=1.5em]
    \item \texttt{snapshot\_to\_json()} performs deterministic JSON serialization
    \item Commit payload is \texttt{serialized\_snapshot || parent\_hash}
    \item SHA-256 hash generated using \texttt{hashlib.sha256(...).hexdigest()}
    \item Parent linkage stored in \texttt{commits.parent\_hash}
    \item Linear-chain constraint enforced by unique key on \texttt{parent\_hash}
\end{itemize}

\textbf{Mathematical form:}
\[
H_i = \mathrm{SHA256}(S_i \Vert H_{i-1})
\]
where $S_i$ is snapshot serialization and $H_{i-1}$ is parent hash.

\subsection{Feature 5: Schema Diff Engine}

\textbf{What it does:}
Generates DDL SQL patch statements between two snapshots.

\textbf{How it is implemented:}
\begin{itemize}[leftmargin=1.5em]
    \item Uses \texttt{sqlglot} to parse \texttt{CREATE TABLE} DDL into AST
    \item Extracts column definitions from \texttt{ColumnDef} nodes
    \item Computes column set delta: added, removed, changed
    \item Emits:\\
    \texttt{ADD COLUMN}, \texttt{DROP COLUMN}, \texttt{MODIFY COLUMN}
    \item Emits table-level \texttt{CREATE TABLE} and \texttt{DROP TABLE} where needed
    \item Orders schema SQL with create/alter first, drop last
\end{itemize}

\textbf{Package vs custom:}
AST parsing is package-provided; actual comparison policy and SQL statement
emission strategy are custom implementation.

\subsection{Feature 6: Data Diff Engine}

\textbf{What it does:}
Generates DML SQL statements to transform rows from old snapshot to new snapshot.

\textbf{How it is implemented:}
\begin{itemize}[leftmargin=1.5em]
    \item Uses PK tuple as stable row identity key
    \item Builds hash maps of old/new rows indexed by PK tuple
    \item Computes three disjoint sets: added, removed, common
    \item Emits \texttt{INSERT} for added rows
    \item Emits \texttt{DELETE} for removed rows
    \item Emits \texttt{UPDATE} only when non-PK values changed
    \item Handles NULL, boolean, numeric, and escaped string literals
\end{itemize}

\textbf{Complexity insight:}
Using map indexing by PK gives approximately linear behavior in number of rows
for each table (average-case $O(n)$ operations for lookup/delta generation).

\textbf{Package vs custom:}
This is fully custom algorithmic code in \texttt{engine/diff.py}.

\subsection{Feature 7: Two-Phase Checkout with Recovery}

\textbf{What it does:}
Safely checks out a target commit into live DB state while handling MySQL DDL
implicit commit behavior.

\textbf{How it is implemented:}
\begin{enumerate}[leftmargin=1.5em]
    \item \textbf{Phase 1 (schema):} execute DDL statements one by one
    \item On schema failure: attempt schema restoration from old snapshot and write
    recovery SQL file to \texttt{.gitdb/recovery\_timestamp.sql}
    \item \textbf{Phase 2 (data):} start transaction, disable FK checks, apply DML,
    re-enable FK checks, commit
    \item On data failure: rollback transaction and raise checkout data error
\end{enumerate}

\begin{tcolorbox}[warningbox]
\textbf{Important MySQL behavior:} many DDL operations auto-commit, so they cannot be
rolled back like normal DML. This is why the schema phase and data phase are
separated in GitDB.
\end{tcolorbox}

\subsection{Feature 8: Status Computation}

\textbf{What it does:}
Compares live database state against HEAD snapshot and reports:
\begin{itemize}[leftmargin=1.5em]
    \item added tables
    \item dropped tables
    \item modified tables (raw DDL changed)
    \item row count deltas per common table
\end{itemize}

\textbf{How it is implemented:}
Custom comparison logic in both CLI and API layers using snapshots from live DB
and current HEAD hash.

\subsection{Feature 9: Command Line Interface}

\textbf{Commands implemented:}
\begin{itemize}[leftmargin=1.5em]
    \item \texttt{gitdb register}
    \item \texttt{gitdb init}
    \item \texttt{gitdb commit}
    \item \texttt{gitdb log} (with \texttt{--oneline} and \texttt{--graph})
    \item \texttt{gitdb diff}
    \item \texttt{gitdb checkout}
    \item \texttt{gitdb status}
\end{itemize}

\textbf{Package vs custom:}
\texttt{click} handles parsing and command registration; \texttt{rich} handles
visual formatting; business behavior is custom orchestration code.

\subsection{Feature 10: REST API for UI Integration}

\textbf{Endpoints:}
\begin{itemize}[leftmargin=1.5em]
    \item \texttt{GET /commits}
    \item \texttt{GET /diff/<hash1>/<hash2>}
    \item \texttt{POST /diff/<hash1>/<hash2>}
    \item \texttt{GET /snapshot/<commit\_hash>}
    \item \texttt{POST /checkout/<commit\_hash>}
    \item \texttt{GET /status}
\end{itemize}

\textbf{How implemented:}
Flask routes in \texttt{api/app.py} call engine modules, serialize response payloads,
and return JSON consumable by React UI.

\subsection{Feature 11: Web Interface}

\textbf{Main UI pages:}
\begin{itemize}[leftmargin=1.5em]
    \item Commit Graph page
    \item Diff Viewer page
    \item Schema Browser page
\end{itemize}

\textbf{Implementation split:}
\begin{itemize}[leftmargin=1.5em]
    \item Custom page and state orchestration in \texttt{ui/src/App.jsx}
    \item Graph rendering via \texttt{@xyflow/react}
    \item Automatic node layout via \texttt{dagre}
    \item Visual split diff via \texttt{@git-diff-view/react}
    \item API access through custom wrappers in \texttt{ui/src/api.js}
    \item Theme system and responsive styling through custom CSS variables and Tailwind
\end{itemize}

\subsection{Feature 12: Automated Testing Strategy}

\textbf{Unit tests:}
\begin{itemize}[leftmargin=1.5em]
    \item \texttt{test\_hash.py}: deterministic hash and Merkle-chain behavior
    \item \texttt{test\_snapshot.py}: snapshot behavior, row-limit guard, no-PK handling
    \item \texttt{test\_diff.py}: schema and data SQL generation
    \item \texttt{test\_checkout.py}: two-phase checkout success/failure paths
    \item \texttt{test\_api.py}: Flask endpoint behavior with monkeypatched DB calls
\end{itemize}

\textbf{Integration tests (live MySQL):}
\begin{itemize}[leftmargin=1.5em]
    \item full round-trip checkout restoration
    \item row-limit error path
    \item schema-only restoration behavior
    \item no-PK table behavior
    \item foreign-key-sensitive checkout flow
\end{itemize}

\section{Data and Control Flow Across Layers}

\subsection{Commit Flow}

\begin{enumerate}[leftmargin=1.5em]
    \item User runs \texttt{gitdb commit -m ... --author ...}
    \item CLI loads repo config and credentials via keyring
    \item Snapshot engine captures target DB state
    \item Hash engine computes commit hash from snapshot plus parent hash
    \item CLI writes commit row and per-table snapshot rows into \texttt{gitdb\_meta}
    \item Repository \texttt{current\_hash} (HEAD) is updated to new hash
\end{enumerate}

\subsection{Checkout Flow}

\begin{enumerate}[leftmargin=1.5em]
    \item User selects target commit (CLI command or UI button via API)
    \item Engine loads old (HEAD) and new snapshots
    \item Diff engine builds schema and data SQL patches
    \item Checkout engine executes schema phase, then transactional data phase
    \item On success, HEAD pointer is updated
\end{enumerate}

\section{What is Built In-House vs Outsourced to Libraries}

\subsection{Primarily In-House (Custom Engineering)}

\begin{itemize}[leftmargin=1.5em]
    \item snapshot model and capture orchestration
    \item commit hash payload design and chain semantics
    \item schema-delta and data-delta generation policies
    \item two-phase checkout safety model
    \item CLI command semantics and user workflows
    \item REST payload design for UI consumption
    \item UI page structure and interaction flow
\end{itemize}

\subsection{Library-Assisted (Infrastructure Building Blocks)}

\begin{itemize}[leftmargin=1.5em]
    \item DB connectivity: \texttt{mysql-connector-python}
    \item DDL parser: \texttt{sqlglot}
    \item secure password hashing: \texttt{argon2-cffi}
    \item secret storage: \texttt{keyring}
    \item CLI framework and formatting: \texttt{click}, \texttt{rich}
    \item API framework: \texttt{flask}, \texttt{flask-cors}
    \item graph/diff frontend widgets: \texttt{@xyflow/react}, \texttt{dagre}, \texttt{@git-diff-view/react}
\end{itemize}

\begin{tcolorbox}[fixbox]
\textbf{Engineering takeaway:}
GitDB is not a wrapper around one ready-made version-control package. Instead, it
combines multiple focused libraries while implementing the core database-versioning
algorithms and workflows as original project code.
\end{tcolorbox}

\section{Current Constraints and Known Limits}

\begin{itemize}[leftmargin=1.5em]
    \item project enforces linear history (no branching support)
    \item row snapshot guard defaults to 10,000 rows per table
    \item tables without primary keys are captured but skipped in row-level diffing
    \item schema rename inference is not semantic; altered definitions can be warned and
    represented through \texttt{MODIFY COLUMN}
    \item recovery file is best-effort schema escape hatch for manual intervention
\end{itemize}

\section{Conclusion}

GitDB uses a layered full-stack architecture and balances reliability with project
scope. The stack selection is pragmatic: mature libraries are used for parsing,
crypto, UI widgets, and connectivity, while the core value of the project
(snapshotting, diff generation, and safe checkout workflow) is implemented through
custom code in the engine layer.

From an academic and engineering perspective, the project demonstrates:
\begin{itemize}[leftmargin=1.5em]
    \item database introspection and metadata modeling
    \item deterministic state serialization and cryptographic identity
    \item SQL patch generation from structural and row-level deltas
    \item safe execution strategy around MySQL DDL transaction limitations
    \item end-to-end productization via CLI, API, UI, and automated tests
\end{itemize}

\end{document}
